\scanpage{257}
\[
\varinjlim\unsure{\varphi_N|\rho_N}\tau_NP
 \longrightarrow\unsure{\varphi_\infty|\rho_\infty}\tau_\infty P
\]
bijectif?
\deleted{\(\operatorname{Hom}\)}
\[\Updownarrow\]
\(\forall M\) de prés fini (il suffit \(M=S[i]\))
\diagram{\begin{array}{ccc}
\operatorname{Hom}_{\mathcal M}(M,\varinjlim\unsure{\varphi_N|\rho_N}\tau_NP)
 &\unsure{-|\longrightarrow}&
\operatorname{Hom}_{\mathcal M}(M,\unsure{\varphi_\infty|\rho_\infty}\tau_\infty P)
 \quad\text{bij.}\\
\simeq&&\simeq\\
\varinjlim\operatorname{Hom}_{\mathcal M}(M,\unsure{\varphi_N|\rho_N}\tau_NP)
 &&\operatorname{Hom}_{\mathcal M_\infty}(\tau_\infty M,\tau_\infty P)\\
\simeq&&\simeq\\
\varinjlim\operatorname{Hom}_{\mathcal M_N}(\tau_NM,\tau_NP)
 &&\underset{\substack{M'\subset M\\M/M'\ \text{négl.}}}{\varinjlim}
 \operatorname{Hom}(M',P)\\
\simeq&&\nearrow\\
\varinjlim\operatorname{Hom}_{\mathcal M}(\tau_NM,P)&&
\end{array}}
dire que (pour \(M\) fixé) c'est un isom pour \(\forall P\), signifie que les deux
proobjets de \(\mathcal M\) formés; soit des \(\tau_NM\), soit des
\(M'\subset M\) tels que \(M/M'\) est négligeable, \deleted{sont} isom par
l'hom can., ou encore que, ces \unsure{ens} de sous-objets sont cofinaux,
i.e que tout \(M'\) contient un \(\tau_NM\)
\inserted{i.e \(M/M'\) est \(\infty\)-négligeable.}
Mais c'est le cas si \(M\) de \deleted{\unsure{p.f.}}\inserted{t.f.}
\deleted{\uline{Hyp} Tout idéal gradué \(J\) de \(S\), tq que \(S/J\) soit
négligeable \unsure{finie.|fini.}}
Donc
\[
\varinjlim\unsure{\varphi_N|\rho_N}\tau_NP
 \overset\sim\longrightarrow\unsure{\varphi_\infty|\rho_\infty}\tau_\infty P
\]
\margin{256}
